\chapter{DMA, PWM, Audio, Pulse and Analog Peripherals}
\label{ch:hal2}

These drivers follow the \emph{claim-a-handle} pattern: a peripheral with several
identical units hands out a limited, controlled handle through a protected pool.
You \texttt{Claim} a unit, use it through the handle, and it is reclaimed when the
handle leaves scope.

\section{GDMA: general DMA}\index{GDMA driver}

\begin{lstlisting}
procedure Claim (C : in out Channel; Peri : Peripheral);
procedure Copy  (C : Channel; Dst, Src : System.Address; Length : Natural);
procedure Start (C : Channel; Dir : Direction; ...);  --  peripheral transfers
function  Done  (C : Channel; Dir : Direction) return Boolean;
procedure Wait  (C : Channel; Dir : Direction);
--  Gapless double-buffered streaming, one ring per direction (see below).
procedure Start_Stream    (C : Channel; Buffer : System.Address; Half_Length : Natural);
function  Await_Half      (C : Channel) return Natural;   --  half to refill (0/1)
procedure Start_In_Stream (C : Channel; Buffer : System.Address; Half_Length : Natural);
function  Await_In_Half   (C : Channel) return Natural;   --  half to read (0/1)
procedure Stop   (C : Channel; Dir : Direction);
procedure Unbind (C : Channel; Dir : Direction);  --  detach an unused direction
\end{lstlisting}

\textbf{Parameters.}
\begin{description}[style=nextline]
\item[\texttt{Peri}] which peripheral the channel serves: \texttt{Mem2Mem} (a
  plain memory copy) or a real peripheral: \texttt{SPI2}, \texttt{SPI3},
  \texttt{UHCI0}, \texttt{I2S0}, \texttt{I2S1}, \texttt{LCD\_CAM}, \texttt{AES},
  \texttt{SHA}, \texttt{ADC\_DAC}, \texttt{RMT}. The pool has five channels.
\item[\texttt{Dir}] transfer direction: \texttt{Mem\_To\_Periph} or
  \texttt{Periph\_To\_Mem}.
\item[\texttt{Dst}, \texttt{Src}, \texttt{Length}] destination/source
  \texttt{'Address}es and a byte count, \texttt{1\,..\,4095}.
\end{description}

\lstexample{Memory-to-memory copy.}
\begin{lstlisting}
with Interfaces;
with ESP32S3.GDMA; use ESP32S3.GDMA;
procedure GDMA_Copy is
   C   : Channel;
   Src : aliased array (0 .. 255) of Interfaces.Unsigned_8 := (others => 16#5A#);
   Dst : aliased array (0 .. 255) of Interfaces.Unsigned_8 := (others => 0);
begin
   Claim (C, Peri => Mem2Mem);                --  a free channel for mem2mem
   Copy  (C, Dst'Address, Src'Address, 256);  --  blocking DMA copy; Dst = Src
end GDMA_Copy;                                 --  C is released on scope exit
\end{lstlisting}

\textbf{Example 2: the RAII guarantee.} Claiming all five channels and then
one more raises (the pool is exhausted); each goes back automatically:
\begin{lstlisting}
with ESP32S3.GDMA; use ESP32S3.GDMA;
procedure GDMA_RAII is
   A, B, C2, D, E : Channel;
begin
   Claim (A, Mem2Mem); Claim (B, Mem2Mem); Claim (C2, Mem2Mem);
   Claim (D, Mem2Mem); Claim (E, Mem2Mem);          --  all five claimed
end GDMA_RAII;                                       --  all five reclaimed here
\end{lstlisting}
(The SPI and I2S drivers use the peripheral \texttt{Start}/\texttt{Wait} path
internally, so most applications use GDMA only for memory copies.)

\texttt{Wait} (and \texttt{Copy}) is interrupt-driven: after a brief spin it
suspends the calling task, and the channel's end-of-transfer interrupt wakes it,
so a transfer hands the core back rather than busy-waiting. The driver owns the
\texttt{Device\_L2\_0} interrupt (CPU\_INT 19) for this.

That suspend is \emph{bounded by a deadline}. A stuck or misconfigured transfer
whose end-of-transfer interrupt never arrives would otherwise block the waiting task
forever, so before it suspends, \texttt{Wait} arms a per-channel
\texttt{Ada.Real\_Time.Timing\_Events} timer a few seconds out; whichever fires
first --- the real EOF interrupt or the timer --- releases the \emph{same}
suspension object, and \texttt{Wait} cancels the other on wake. \texttt{Set\_True}
is idempotent, so the two racing near the deadline is harmless. The normal path
keeps its exact, zero-latency interrupt wake (the timer is cancelled before it ever
runs); only a genuine fault trips the deadline, and then \texttt{Wait} returns
rather than deadlocking the system. A timed entry call would read more directly, but
the Ravenscar/Jorvik runtime forbids \texttt{select}, so a timing-event wake is the
profile-legal way to put a ceiling on a \texttt{Suspend\_Until\_True}.

\textbf{Streaming rings.} \texttt{Start\_Stream} services a continuous
\emph{producer}: it links two descriptors, one per half of the caller's buffer,
into a ring the OUT engine walks forever, and enables the per-descriptor
completion interrupt. \texttt{Await\_Half} suspends until a half has drained
and returns which one (0 or 1) is now safe to refill. The refill is therefore
paced by the DMA itself, not by a timer, so live-generated audio never drifts
against the peripheral clock. \texttt{Start\_In\_Stream} and
\texttt{Await\_In\_Half} are the receive mirror for a continuous
\emph{consumer} (a demodulator, a recorder): the IN engine fills the two halves
forever, and the caller reads each half as it completes, missing no samples.
Run one ring in each direction on two channels and you have true full duplex.
\texttt{Stop} with the matching direction ends a ring.

\textbf{One channel per direction.} \texttt{Claim} binds \emph{both}
directions of a channel to the peripheral, but the hardware routes each
peripheral request line to only one channel per direction. When two channels
serve the same peripheral (one streaming TX, another streaming RX), the idle
binding on one channel \emph{shadows} the active binding on the other, and that
direction silently transfers nothing. (Measured: a capture stream beside a
running I2S transmit stream delivered zero bytes until the transmit channel's
unused IN binding was detached.) \texttt{Unbind} detaches one direction; every
claimant that uses a channel one-way should unbind the other, and the I2S driver
now does.

\section{MCPWM: motor-control PWM}\index{MCPWM driver}

\begin{lstlisting}
--  The unit's clock comes up on the first Claim of any of its channels.
procedure Claim (C : in out Channel; Unit : MCPWM_Unit; Index : Channel_Index);
procedure Configure_Channel (C : Channel; Freq : Positive; Pin : Pin_Id;
                             Complement_Pin : Optional_Pin := No_Pin;
                             Dead_Time_Ns   : Natural := 0);
procedure Start (C : Channel);   procedure Set_Duty (C : Channel; Percent : Duty_Percent);
\end{lstlisting}

\textbf{Parameters.}
\begin{description}[style=nextline]
\item[\texttt{Unit}] \texttt{MCPWM0} or \texttt{MCPWM1}.
\item[\texttt{Index}] generator channel within the unit: \texttt{Ch0},
  \texttt{Ch1} or \texttt{Ch2}.
\item[\texttt{Freq}] PWM frequency in Hz.
\item[\texttt{Pin}] the output pad.
\item[\texttt{Complement\_Pin}] optional inverted output for a half-bridge
  (\texttt{Optional\_Pin}, default \texttt{No\_Pin}).
\item[\texttt{Dead\_Time\_Ns}] gap inserted between the complementary edges, in
  nanoseconds; default \texttt{0}.
\item[\texttt{Percent}] duty cycle, \texttt{0.0\,..\,100.0}.
\end{description}

\lstexample{20\,kHz at 25\,\% on one pin.}
\begin{lstlisting}
with ESP32S3.MCPWM; use ESP32S3.MCPWM;
procedure MCPWM_Simple is
   C : Channel;
begin
   Claim (C, MCPWM0, Ch0);   --  first Claim brings up the MCPWM0 clock
   Configure_Channel (C, Freq => 20_000, Pin => 4);
   Set_Duty (C, 25.0);
   Start (C);
end MCPWM_Simple;
\end{lstlisting}

\lstexample{A complementary pair with 500\,ns dead-time (half-bridge).}
\begin{lstlisting}
with ESP32S3.MCPWM; use ESP32S3.MCPWM;
procedure MCPWM_Pair is
   C : Channel;
begin
   Claim (C, MCPWM0, Ch0);   --  first Claim brings up the MCPWM0 clock
   Configure_Channel (C, Freq => 20_000, Pin => 4,
                      Complement_Pin => 5, Dead_Time_Ns => 500);
   Set_Duty (C, 60.0);
   Start (C);
end MCPWM_Pair;
\end{lstlisting}
The driver also offers a carrier (\texttt{Set\_Carrier}: \texttt{Prescale}
\texttt{0\,..\,15}, \texttt{Duty\_Eighths} \texttt{1\,..\,7}, \texttt{First\_Pulse}
\texttt{0\,..\,15}) and fault/trip-zone protection: \texttt{Configure\_Fault}
binds a \texttt{Fault\_Input} (\texttt{Fault0\,..\,2}) to a pad, and
\texttt{Protect\_Channel} sets the \texttt{Fault\_Mode} (\texttt{One\_Shot} or
\texttt{Cycle\_By\_Cycle}) and \texttt{Trip\_Action} (\texttt{No\_Change},
\texttt{Force\_Low} or \texttt{Force\_High}).

\section{I2S: digital audio over GDMA}\index{I2S driver}

\begin{lstlisting}
type I2S_Mode is (Standard, PDM);
type PCM_8  is array (Natural range <>) of Interfaces.Integer_8;   -- the typed,
type PCM_16 is array (Natural range <>) of Interfaces.Integer_16;  --  signed PCM
type PCM_32 is array (Natural range <>) of Interfaces.Integer_32;  --  buffers
--  Acquire claims a port AND brings it up (audio format + pin routing).  The
--  first Acquire of a port opens it; a later Acquire reuses it (Reconfigure
--  re-opens to change the format), so a codec's MCLK stays running.
procedure Acquire (S : in out Session; Port : I2S_Port;
                   Sample_Rate : Positive := 16_000; Bits : Sample_Bits := Bits_16;
                   Mode : I2S_Mode := Standard;
                   Bclk, Ws, Dout, Din, Mclk : Optional_Pin := No_Pin);
procedure Reconfigure (S : Session; ...);  -- re-open the held port's format
function  Configured_Bits (S : Session) return Sample_Bits;
--  Typed transfers (preferred): the element type fixes the on-wire width, so the
--  byte count is derived and a precondition checks it against the port's Bits.
procedure Write    (S : Session; Samples : PCM_16)
   with Pre => Configured_Bits (S) = Bits_16;
procedure Read     (S : Session; Samples : out PCM_16)
   with Pre => Configured_Bits (S) = Bits_16;
procedure Transfer (S : Session; Tx : PCM_16; Rx : out PCM_16)
   with Pre => Configured_Bits (S) = Bits_16 and then Tx'Length = Rx'Length;
--  PCM_8 / PCM_32 overloads likewise (also Start_Continuous / Capture).  The
--  *_Raw forms take a 'Address + byte Length for bytes already framed elsewhere:
procedure Write_Raw (S : Session; Tx : System.Address; Length : Natural);
--  Gapless double-buffered streaming, both directions (see below).
procedure Start_Stream (S : Session; Samples : PCM_16);   --  TX ring, 2 halves
function  Await_Half   (S : Session) return Natural;      --  half to refill (0/1)
procedure Start_Capture_Stream (S : Session; Rx : System.Address;
                                Half_Length : Natural);   --  RX ring, own channel
function  Await_Capture_Half   (S : Session) return Natural;  --  half to read
procedure Stop_Capture         (S : Session);
\end{lstlisting}

\textbf{Parameters.}
\begin{description}[style=nextline]
\item[\texttt{Port}] \texttt{I2S0} or \texttt{I2S1}.
\item[\texttt{Sample\_Rate}] frames per second (\texttt{Positive}); default
  \texttt{16\_000}.
\item[\texttt{Bits}] sample width: \texttt{Bits\_8}, \texttt{Bits\_16} (default),
  \texttt{Bits\_24} or \texttt{Bits\_32}.
\item[\texttt{Mode}] data path: \texttt{Standard} (default, plain I2S/TDM) or
  \texttt{PDM} (the sigma-delta converters; see below).
\item[\texttt{Bclk}, \texttt{Ws}, \texttt{Dout}, \texttt{Din}] the bit-clock,
  word-select, data-out and data-in lines (\texttt{Optional\_Pin};
  \texttt{Dout}/\texttt{Din} default \texttt{No\_Pin}).
\item[\texttt{Samples}] a typed \texttt{PCM\_8\,/\,16\,/\,32} buffer; the driver
  derives the byte count and a precondition checks the width against the port.
  The \texttt{*\_Raw} forms instead take a \texttt{'Address} + byte
  \texttt{Length} (\texttt{1\,..\,4095}) for bytes already framed elsewhere.
\end{description}

The element type of a \texttt{PCM\_*} buffer fixes the on-wire width, so a
\texttt{Write\,(S,\,Samples)} needs no caller-side \texttt{'Address} or
\texttt{*\,2}: the driver computes the byte count from \texttt{Samples'Length},
and the precondition \texttt{Configured\_Bits\,(S)\,=\,Bits\_16} turns a
buffer-vs.-port width mismatch from silent noise into a caught contract
violation. Reach for \texttt{Write\_Raw}/\texttt{Transfer\_Raw} only for bytes
that are \emph{already} framed (a DMA buffer built elsewhere, or an opaque
bit-pattern self-test).

A one-shot \texttt{Write} also has to know when it is truly \emph{finished}. The
GDMA completing means only that the samples have reached the I2S TX FIFO --- not
that the serializer has clocked them out onto the wire. Clearing \texttt{TX\_START}
the instant the DMA is done would discard whatever is still in the FIFO, truncating
the tail of every buffer (the end of the sound simply cut off). So after the DMA
wait the driver spins on the \texttt{STATE.TX\_IDLE} status --- FIFO drained,
serializer idle --- before it stops the channel, so the last samples actually leave
the pin.

\lstexample{Play a buffer of samples.}
\begin{lstlisting}
with ESP32S3.I2S; use ESP32S3.I2S;
procedure I2S_Play is
   S    : Session;
   Tone : PCM_16 (0 .. 255) := (others => 0);
begin
   Acquire (S, I2S0, Sample_Rate => 44_100, Bits => Bits_16,
            Bclk => 4, Ws => 5, Dout => 6);   --  claim + configure
   Write (S, Tone);                           --  256 samples x 2 bytes, derived
end I2S_Play;
\end{lstlisting}

\lstexample{Wiring-free full-duplex loopback self-test.}
\begin{lstlisting}
with ESP32S3.I2S; use ESP32S3.I2S;
procedure I2S_Loopback is
   S  : Session;
   Tx : PCM_16 (0 .. 63) := (others => 0);    --  fill with a test pattern
   Rx : PCM_16 (0 .. 63) := (others => 0);
begin
   Acquire (S, I2S0);                          --  claim (defaults: 16k 16-bit)
   Enable_Loopback (S, Pad => 7);              --  TX folded to RX (held port)
   Transfer (S, Tx, Rx);                       --  Rx = Tx (same length checked)
end I2S_Loopback;
\end{lstlisting}

\subsection{Gapless streaming and full duplex}

A one-shot \texttt{Write} stops the clock when it finishes, and back-to-back
writes leave an audible gap. \texttt{Start\_Continuous} closes the gap for a
\emph{static} buffer (it replays one buffer forever), but live-generated audio
needs to refill while playing, and pacing the refill from a CPU timer drifts
against the I2S clock. (Measured: a timer-paced double buffer feeding a codec
produced a beat that returned every few seconds as the two clocks slid past each
other.) \texttt{Start\_Stream} fixes the pacing at the source: the DMA loops
the two halves of the buffer forever, and \texttt{Await\_Half} blocks until a
half has drained, returning which one (0 or 1) to refill. The producer is paced
by the same clock that consumes the samples, so it can never drift.

\texttt{Start\_Capture\_Stream} is the receive mirror. It claims its
\emph{own} DMA channel, so it runs concurrently with a streaming transmit:
play out of one ring while recording into the other, both gapless, with no
clock interruption in either direction (true full duplex).
\texttt{Await\_Capture\_Half} returns each captured half as the DMA fills
it, so a continuous consumer (a recorder, a software modem) misses no samples;
\texttt{Stop\_Capture} ends it.

Two facts make external capture work, and both are invisible until they bite.
First, the RX block runs as a \emph{slave}: it samples on the bit and word
clocks it receives, not on clocks it generates. The driver therefore feeds the
RX clock inputs from the very pads the TX master drives, so capture samples on
the true on-wire clock, phase-exact with the external device. (Without that
routing, slave RX has no clock at all and capture waits forever; the internal
loopback self-test masks the bug, because \texttt{SIG\_LOOPBACK} bypasses the
pads.) Second, the one-channel-per-direction rule from the GDMA section: the
transmit stream's channel unbinds its unused RX direction and the capture
channel unbinds its unused TX direction, or one shadows the other.

\lstexample{Full duplex: synthesize live audio while recording the input.}
\begin{lstlisting}
Half : constant := 500;                        --  31 ms per half at 16 kHz
Play_Ring, Capture_Ring : PCM_16 (0 .. 4 * Half - 1);
S : Session;
begin
   Acquire (S, I2S0, Bclk => 4, Ws => 5, Dout => 6, Din => 7);
   Fill (Play_Ring);                           --  prime both halves
   Start_Stream (S, Play_Ring);
   Start_Capture_Stream (S, Capture_Ring'Address, 2 * Half * 2);
   loop
      Refill  (Play_Ring,    Await_Half (S));          --  DMA-paced, no drift
      Consume (Capture_Ring, Await_Capture_Half (S));  --  misses no samples
   end loop;
\end{lstlisting}

\subsection{PDM mode: the hardware PCM\,$\leftrightarrow$\,PDM converters}

\texttt{Mode => PDM} routes the data path through the sigma-delta converters
instead of the plain serial framer. PDM (pulse-density modulation) is the
sigma-delta sibling of PWM: a 1-bit oversampled stream whose \emph{density} of
ones encodes the analog level. The S3 I2S converts in both directions, and the
DMA still moves ordinary PCM, so only the on-wire format changes:

\begin{itemize}
  \item \textbf{TX} (\texttt{PCM2PDM}) turns each PCM sample into a 1-bit stream:
        feed it to a class-D amplifier or an RC low-pass for analog audio out
        (the S3 has no analog DAC, so this is how you get one).
  \item \textbf{RX} (\texttt{PDM2PCM}) decimates a 1-bit PDM input back to PCM.
        PDM is the format \emph{emitted} by digital \textbf{PDM microphones} and
        some PDM-output \textbf{ADCs}; the converter turns their bitstream into the
        PCM samples your program reads. In PDM the receiver is the clock master:
        it drives the device clock out on \texttt{Ws} and reads the bitstream in on
        \texttt{Din}.
\end{itemize}

\begin{lstlisting}
with ESP32S3.I2S; use ESP32S3.I2S;
procedure PDM_Mic is
   S  : Session;
   Rx : PCM_16 (0 .. 511) := (others => 0);
begin
   Acquire (S, I2S0, Sample_Rate => 16_000, Bits => Bits_16, Mode => PDM,
            Ws => 5, Din => 6);               --  claim + configure (PDM mic)
   Read (S, Rx);                              --  PDM mic -> signed PCM_16
end PDM_Mic;
\end{lstlisting}

\textbf{A note on verifying PDM.} Every other HAL driver self-tests with no
external hardware; PDM is the exception, and the reason is instructive. There is
no on-chip loopback for the converters (the \texttt{SIG\_LOOPBACK} bit only
shares the \emph{standard}-I2S WS and BCK, not the bit timing between
\texttt{PCM2PDM} and \texttt{PDM2PCM}), so an internal
\texttt{PCM}\,$\rightarrow$\,\texttt{PDM}\,$\rightarrow$\,\texttt{PCM} loop
recovers nothing. Nor can a static level
synthesised from a GPIO (a pull resistor or a driven pin) stand in for a mic: the
decimator's mandatory high-pass filter strips DC, so a constant 1 or 0 is rejected.
Genuine PDM verification needs a real device producing a toggling bitstream, so
the \texttt{esp32s3\_i2s\_pdm} example is a microphone-capture demo rather than a
self-test.

\section{LEDC: LED / general PWM}\index{LEDC driver}

\begin{lstlisting}
procedure Claim (C : in out Channel; Index : Channel_Index);
procedure Configure (C : in out Channel; Freq : Positive; Pin : Pin_Id;
                     Bits : Resolution := 10);
procedure Set_Duty (C : Channel; Percent : Duty_Percent);
\end{lstlisting}

\textbf{Parameters.}
\begin{description}[style=nextline]
\item[\texttt{Index}] channel \texttt{0\,..\,7} (\texttt{Channel\_Index}).
\item[\texttt{Freq}] PWM frequency in Hz.
\item[\texttt{Pin}] the output pad.
\item[\texttt{Bits}] duty resolution in bits, \texttt{1\,..\,14}
  (\texttt{Resolution}); default \texttt{10}. More bits give finer duty steps but a
  lower maximum frequency.
\item[\texttt{Percent}] duty cycle, \texttt{0.0\,..\,100.0}.
\end{description}

\lstexample{5\,kHz at 50\,\%.}
\begin{lstlisting}
with ESP32S3.LEDC; use ESP32S3.LEDC;
procedure LEDC_Simple is
   C : Channel;
begin
   Claim (C, 0);
   Configure (C, Freq => 5_000, Pin => 2);
   Set_Duty (C, 50.0);
end LEDC_Simple;
\end{lstlisting}

\lstexample{Fade an LED.}
\begin{lstlisting}
with Ada.Real_Time; use Ada.Real_Time;
with ESP32S3.LEDC;   use ESP32S3.LEDC;
procedure LEDC_Fade is
   C : Channel;
begin
   Claim (C, 0);
   Configure (C, Freq => 5_000, Pin => 2);
   for D in 0 .. 100 loop
      Set_Duty (C, Duty_Percent (D));
      delay until Clock + Milliseconds (10);
   end loop;
end LEDC_Fade;
\end{lstlisting}

\section{RMT: remote-control pulse TX/RX (IR, WS2812)}\index{RMT driver}

\begin{lstlisting}
type RMT_Symbol is record
   Level0, Level1 : Boolean := False; Duration0, Duration1 : Tick_Count := 0;
end record;
procedure Configure (C : in out TX_Channel; Resolution_Hz : Positive; Pin : Pin_Id);
procedure Transmit  (C : TX_Channel; Symbols : Symbol_Array);
procedure Configure (C : in out RX_Channel; Resolution_Hz : Positive; Pin : Pin_Id);
procedure Start (C : RX_Channel);
procedure Receive (C : RX_Channel; Into : out Symbol_Array; Count : out Natural);
\end{lstlisting}

\textbf{Parameters.}
\begin{description}[style=nextline]
\item[\texttt{Index}] channel \texttt{0\,..\,3}: four TX channels
  (\texttt{TX\_Index}) and four RX (\texttt{RX\_Index}), separate pools.
\item[\texttt{Resolution\_Hz}] the tick clock: one \texttt{Tick\_Count} unit lasts
  \texttt{1/Resolution\_Hz} seconds.
\item[\texttt{Pin}] the carrier pad.
\item[\texttt{Idle\_Ticks}] (RX) the silence that ends a frame,
  \texttt{0\,..\,32\_767}; default \texttt{2\_000}.
\item[\texttt{Symbols} / \texttt{Into}] a \texttt{Symbol\_Array} of
  \texttt{RMT\_Symbol}s, each holding two level+duration halves;
  \texttt{Duration0}/\allowbreak\texttt{Duration1} are \texttt{Tick\_Count}
  (\texttt{0\,..\,32\_767}). One symbol-RAM block holds 48 symbols; a longer burst
  is streamed by refilling the block in halves (wrap mode).
\item[\texttt{Count}] (\texttt{Receive}) \texttt{out} count of symbols captured.
\end{description}

\lstexample{Transmit a pulse burst (1\,\textmu s resolution).}
\begin{lstlisting}
with ESP32S3.RMT; use ESP32S3.RMT;
procedure RMT_Tx is
   Tx    : TX_Channel;
   Burst : constant Symbol_Array :=
     ((Level0 => True,  Duration0 => 10, Level1 => False, Duration1 => 10),
      (Level0 => True,  Duration0 => 20, Level1 => False, Duration1 => 20));
begin
   Claim (Tx, 0);
   Configure (Tx, Resolution_Hz => 1_000_000, Pin => 4);
   Transmit (Tx, Burst);
end RMT_Tx;
\end{lstlisting}

\lstexample{Receive into a buffer (e.g. an IR remote).}
\begin{lstlisting}
with ESP32S3.RMT; use ESP32S3.RMT;
procedure RMT_Rx is
   Rx   : RX_Channel;
   Syms : Symbol_Array (0 .. 63);
   N    : Natural;
begin
   Claim (Rx, 0);
   Configure (Rx, Resolution_Hz => 1_000_000, Pin => 5);
   Start (Rx);
   Receive (Rx, Syms, N);   --  Syms (0 .. N-1) are the captured pulses
end RMT_Rx;
\end{lstlisting}

\section{PCNT: pulse / edge counter}\index{PCNT driver}

\begin{lstlisting}
procedure Configure (U : in out Unit; Pin : Pin_Id; Both_Edges : Boolean := False);
function  Count (U : Unit) return Integer;   procedure Clear (U : Unit);
\end{lstlisting}

\textbf{Parameters.}
\begin{description}[style=nextline]
\item[\texttt{Index}] unit \texttt{0\,..\,3} (\texttt{Unit\_Index}).
\item[\texttt{Pin}] the input pad whose edges are counted.
\item[\texttt{Both\_Edges}] \texttt{False} (default) counts rising edges only;
  \texttt{True} counts both (double resolution).
\item[\texttt{Count}] returns the signed counter (16-bit, wraps at $\pm$32768).
\end{description}

\lstexample{Count rising edges on a pin.}
\begin{lstlisting}
with Ada.Text_IO;  use Ada.Text_IO;
with ESP32S3.PCNT; use ESP32S3.PCNT;
procedure PCNT_Count is
   U : Unit;
begin
   Claim (U, 0);
   Configure (U, Pin => 6);
   Clear (U);
   --  ... pulses arrive ...
   Put_Line ("edges =" & Integer'Image (Count (U)));
end PCNT_Count;
\end{lstlisting}

\lstexample{Count both edges (double the resolution).}
\begin{lstlisting}
with ESP32S3.PCNT; use ESP32S3.PCNT;
procedure PCNT_Both is
   U : Unit;
begin
   Claim (U, 0);
   Configure (U, Pin => 6, Both_Edges => True);
end PCNT_Both;
\end{lstlisting}

\section{SDM: sigma-delta density output}\index{SDM driver}

\begin{lstlisting}
procedure Configure (C : in out Channel; Pin : Pin_Id;
                     Carrier_Hz : Positive := 1_000_000);
procedure Set_Density (C : Channel; Percent : Density_Percent);
\end{lstlisting}

\textbf{Parameters.}
\begin{description}[style=nextline]
\item[\texttt{Index}] channel \texttt{0\,..\,7} (\texttt{Channel\_Index}).
\item[\texttt{Pin}] the output pad.
\item[\texttt{Carrier\_Hz}] the desired pulse-stream (carrier) frequency in Hz;
  default \texttt{1\_000\_000}. The modulator divides the APB clock
  ($\approx$80\,MHz) by an integer \texttt{1\,..\,256}, so the achieved rate is
  the nearest \texttt{APB/N} (roughly \texttt{312\_500\,..\,80\_000\_000}\,Hz),
  rounded for you. A higher carrier sits further above your RC low-pass cutoff and
  smooths to a cleaner level; the exact value rarely matters.
\item[\texttt{Percent}] target density, \texttt{0.0\,..\,100.0}.
\end{description}

\lstexample{A 25\,\% density bitstream (analog-ish level via an RC filter).}
\begin{lstlisting}
with ESP32S3.SDM; use ESP32S3.SDM;
procedure SDM_Density is
   C : Channel;
begin
   Claim (C, 0);
   Configure (C, Pin => 7, Carrier_Hz => 1_000_000);
   Set_Density (C, 25.0);
end SDM_Density;
\end{lstlisting}

\lstexample{Breathe an LED by sweeping density.}
\begin{lstlisting}
with Ada.Real_Time; use Ada.Real_Time;
with ESP32S3.SDM;    use ESP32S3.SDM;
procedure SDM_Breathe is
   C : Channel;
begin
   Claim (C, 0);
   Configure (C, Pin => 7);
   for P in 0 .. 100 loop
      Set_Density (C, Density_Percent (P));
      delay until Clock + Milliseconds (10);
   end loop;
end SDM_Breathe;
\end{lstlisting}

\section{TWAI: CAN 2.0}\index{TWAI driver}

\begin{lstlisting}
subtype Standard_Id is Interfaces.Unsigned_32 range 0 .. 16#7FF#;        -- 11-bit
subtype Extended_Id is Interfaces.Unsigned_32 range 0 .. 16#1FFF_FFFF#;  -- 29-bit
type Standard_Frame is record   -- 11-bit id (CAN 2.0A)
   Id : Standard_Id := 0; Remote : Boolean := False;   -- Remote => RTR request
   Length : Data_Length := 0; Data : Data_Bytes := (others => 0);
end record;
type Extended_Frame is record   -- 29-bit id (CAN 2.0B)
   Id : Extended_Id := 0; Remote : Boolean := False;
   Length : Data_Length := 0; Data : Data_Bytes := (others => 0);
end record;
procedure Acquire (S : in out Session;    -- claim AND configure the controller
                   Mode : Bus_Mode := Normal; Bit_Rate : Positive := 125_000);
procedure Reconfigure (S : Session;       -- re-apply mode + bit rate, held
                       Mode : Bus_Mode := Normal; Bit_Rate : Positive := 125_000);
procedure Send (S : Session; F : Standard_Frame);   -- overloaded on the frame type
procedure Send (S : Session; F : Extended_Frame);
function  Available   (S : Session) return Boolean; -- a frame is waiting
function  Is_Extended (S : Session) return Boolean; -- ... and it is 29-bit
procedure Receive (S : Session; F : out Standard_Frame; Got : out Boolean);
procedure Receive (S : Session; F : out Extended_Frame; Got : out Boolean);
\end{lstlisting}

\textbf{Parameters.}
\begin{description}[style=nextline]
\item[\texttt{Mode}] \texttt{Normal} (default), \texttt{Listen\_Only} (receive
  without ACKing) or \texttt{Self\_Test} (ACK its own frames, the wiring-free
  loopback).
\item[\texttt{Bit\_Rate}] bus speed in bits/s; default \texttt{125\_000} (e.g.
  \texttt{250\_000}, \texttt{500\_000}, \texttt{1\_000\_000}).
\item[\texttt{Tx}, \texttt{Rx}] the bus pads (\texttt{Optional\_Pin}).
\item[\texttt{F} (\texttt{Standard\_Frame} / \texttt{Extended\_Frame})] the frame:
  \texttt{Id} (range-checked to its width: \texttt{Standard\_Id}
  \texttt{0\,..\,16\#7FF\#} or \texttt{Extended\_Id}
  \texttt{0\,..\,16\#1FFF\_FFFF\#}), \texttt{Length} (\texttt{0\,..\,8} bytes),
  \texttt{Data} (an 8-byte array), and \texttt{Remote}: \texttt{False} (default)
  for a data frame, \texttt{True} for a remote-transmission request (RTR), which
  carries \texttt{Id} and \texttt{Length} but no \texttt{Data}. \texttt{Send} is
  overloaded on the frame type.
\item[\texttt{Got}] (\texttt{Receive}) \texttt{out Boolean}: \texttt{True} if a
  frame of the requested width was read.
\item[\texttt{Pad}] (\texttt{Enable\_Loopback}) the single self-test pad.
\end{description}

\textbf{Standard vs.\ extended frames.} A CAN identifier is either 11-bit
(standard) or 29-bit (extended, CAN 2.0B). The two are distinct frame types, so
the id is range-checked to its width and \texttt{Send} is overloaded: you
cannot put a 29-bit id in a standard frame. On receive the \emph{sender} picks
the width, so you peek first: \texttt{Available} reports a waiting frame,
\texttt{Is\_Extended} reports its width, and you call the matching
\texttt{Receive} overload (Example 2). Orthogonal to the width, \texttt{Remote}
marks a \emph{remote-transmission request} (RTR): a frame that carries an id and a
length but no data, asking whichever node owns that id to reply with the data.
Set it on either frame type; a received frame reports it.

\lstexample{Send a standard and an extended frame.}
\begin{lstlisting}
with ESP32S3.TWAI; use ESP32S3.TWAI;
procedure TWAI_Send is
   S : Session;
begin
   Acquire (S, Mode => Normal, Bit_Rate => 500_000);   --  own + configure
   Configure_Pins (S, Tx => 4, Rx => 5);               --  then route the pins
   Send (S, Standard_Frame'(Id => 16#123#, Remote => False, Length => 2,
                            Data => (16#DE#, 16#AD#, others => 0)));
   Send (S, Extended_Frame'(Id => 16#14AB_CDE#, Remote => False, Length => 1,
                            Data => (0 => 16#42#, others => 0)));  -- type picks Send
   Send (S, Standard_Frame'(Id => 16#7A5#, Remote => True,        -- RTR: request
                            Length => 8, Data => (others => 0)));  -- 8 bytes, none sent
end TWAI_Send;
\end{lstlisting}

\lstexample{Self-test loopback (no transceiver, no wiring).}
\begin{lstlisting}
with ESP32S3.TWAI; use ESP32S3.TWAI;
procedure TWAI_Selftest is
   S   : Session;
   RE  : Extended_Frame;
   Got : Boolean := False;
begin
   Acquire (S, Mode => Self_Test);
   Enable_Loopback (S, Pad => 4);
   Send (S, Extended_Frame'(Id => 16#14AB_CDE#, Remote => False, Length => 1,
                            Data => (0 => 16#42#, others => 0)));
   if Available (S) and then Is_Extended (S) then
      Receive (S, RE, Got);   --  Got, RE.Id = 16#14AB_CDE#, RE.Remote = False
   end if;
end TWAI_Selftest;
\end{lstlisting}

\section{Timer: general-purpose timers}\index{Timer driver}

\begin{lstlisting}
procedure Configure (T : in out Timer; Tick_Hz : Positive := 1_000_000);
procedure Start (T : Timer);  function Value (T : Timer) return Ticks;
procedure Set_Alarm (T : Timer; At_Ticks : Ticks);
function  Alarm_Fired (T : Timer) return Boolean;  procedure Clear_Alarm (T : Timer);
\end{lstlisting}

\textbf{Parameters.}
\begin{description}[style=nextline]
\item[\texttt{Index}] \texttt{0} or \texttt{1} (\texttt{Timer\_Index}, the two
  timer groups TIMG0/TIMG1).
\item[\texttt{Tick\_Hz}] counter rate in Hz; default \texttt{1\_000\_000}
  (1 tick = 1\,\textmu s).
\item[\texttt{At\_Ticks} / \texttt{Value}] a 54-bit count (\texttt{Ticks}, an
  \texttt{Unsigned\_64}); \texttt{Set\_Alarm} fires when the counter reaches
  \texttt{At\_Ticks}.
\end{description}

\lstexample{Measure an elapsed interval in microseconds.}
\begin{lstlisting}
with Ada.Text_IO;   use Ada.Text_IO;
with ESP32S3.Timer; use ESP32S3.Timer;
procedure Timer_Measure is
   T       : Timer;
   Start_T : Ticks;
begin
   Claim (T, 0);
   Configure (T, Tick_Hz => 1_000_000);   --  1 tick = 1 us
   Start (T);
   Start_T := Value (T);
   --  ... work ...
   Put_Line ("elapsed us =" & Ticks'Image (Value (T) - Start_T));
end Timer_Measure;
\end{lstlisting}

\lstexample{A one-shot alarm 50\,ms out.}
\begin{lstlisting}
with ESP32S3.Timer; use ESP32S3.Timer;
procedure Timer_Alarm is
   T : Timer;
begin
   Claim (T, 0);
   Configure (T, Tick_Hz => 1_000_000);
   Start (T);
   Set_Alarm (T, Value (T) + 50_000);
   loop exit when Alarm_Fired (T); end loop;
   Clear_Alarm (T);
end Timer_Alarm;
\end{lstlisting}

\section{LCD: 8-bit i80 parallel master}\index{LCD driver}

\begin{lstlisting}
--  Acquire claims the controller AND brings it up (pixel clock + pin routing):
procedure Acquire (S : in out Session; Pclk_Hz : Positive := 1_000_000;
                   Data : Data_Pins := (others => No_Pin);
                   Pclk : Optional_Pin := No_Pin);
procedure Reconfigure (S : Session; ...);  -- re-apply clock + pins, held
procedure Transmit (S : Session; Tx : System.Address; Length : Natural;
                    Ok : out Boolean);
\end{lstlisting}

\textbf{Parameters.}
\begin{description}[style=nextline]
\item[\texttt{Pclk\_Hz}] pixel clock in Hz; default \texttt{1\_000\_000} (the
  achieved value is \texttt{20\,MHz} divided by an integer).
\item[\texttt{Data}] a \texttt{Data\_Pins} value: an 8-element array of
  \texttt{Optional\_Pin} for the data lines D0\,..\,D7.
\item[\texttt{Pclk} / \texttt{Pclk\_Pad}] the pixel-clock pad.
\item[\texttt{Tx}, \texttt{Length}] frame-buffer \texttt{'Address} and a byte
  count, \texttt{1\,..\,4095}.
\item[\texttt{Ok}] \texttt{out Boolean}: \texttt{True} on a completed transfer.
\end{description}

\lstexample{Push a frame buffer to a parallel display.}
\begin{lstlisting}
with Interfaces;
with ESP32S3.LCD; use ESP32S3.LCD;
procedure LCD_Frame is
   S  : Session;
   Ok : Boolean;
   FB : aliased array (0 .. 3999) of Interfaces.Unsigned_8 := (others => 16#FF#);
begin
   Acquire (S, Pclk_Hz => 200_000,
            Data => (1, 2, 3, 4, 5, 6, 7, 8), Pclk => 9);  --  claim + configure
   Transmit (S, FB'Address, FB'Length, Ok);
end LCD_Frame;
\end{lstlisting}

\lstexample{Expose the pixel clock on a pad (for a scope).}
\begin{lstlisting}
with ESP32S3.LCD; use ESP32S3.LCD;
procedure LCD_Clk is
   Lcd : Session;
begin
   Acquire (Lcd, Pclk_Hz => 1_000_000);    -- claim + bring up the clock
   Enable_Clock_Out (Lcd, Pclk_Pad => 9);  -- expose PCLK on IO9
end LCD_Clk;
\end{lstlisting}

\section{ADC: SAR analog input}\index{ADC driver}

\begin{lstlisting}
procedure Claim (R : in out Reader; Unit : ADC_Unit);
function Read (R : Reader; Ch : Channel_Index; Atten : Attenuation := Db_12)
   return Raw_Value;                              --  0 .. 4095 (12-bit)
\end{lstlisting}

\textbf{Parameters.}
\begin{description}[style=nextline]
\item[\texttt{Unit}] \texttt{ADC1} or \texttt{ADC2}.
\item[\texttt{Ch}] channel \texttt{0\,..\,9} (\texttt{Channel\_Index}); each maps to
  a fixed pad.
\item[\texttt{Atten}] input attenuation, trading range for resolution:
  \texttt{Db\_0} ($\sim$1.1\,V full scale), \texttt{Db\_2\_5} ($\sim$1.5\,V),
  \texttt{Db\_6} ($\sim$2.2\,V) or \texttt{Db\_12} ($\sim$3.3\,V, default).
\item[returns] a \texttt{Raw\_Value}, \texttt{0\,..\,4095} (12-bit).
\end{description}

\lstexample{Read a 12-bit sample.}
\begin{lstlisting}
with Ada.Text_IO; use Ada.Text_IO;
with ESP32S3.ADC; use ESP32S3.ADC;
procedure ADC_Read_Sample is
   R : Reader;
   V : Raw_Value;
begin
   Claim (R, ADC1);
   V := Read (R, Ch => 0);     --  0 .. 4095 at full (12 dB) attenuation
   Put_Line (Raw_Value'Image (V));
end ADC_Read_Sample;
\end{lstlisting}

\lstexample{A lower attenuation for a small signal.}
\begin{lstlisting}
with Ada.Text_IO; use Ada.Text_IO;
with ESP32S3.ADC; use ESP32S3.ADC;
procedure ADC_Atten is
   R : Reader;
   V : Raw_Value;
begin
   Claim (R, ADC1);
   V := Read (R, Ch => 3, Atten => Db_0);   --  ~0..1.1 V full scale
   Put_Line (Raw_Value'Image (V));
end ADC_Atten;
\end{lstlisting}

\textbf{Per-attenuation calibration.} The SAR ADC needs a self-calibration
``init code'' --- found once at start-up by grounding the input and binary-searching
the code that reads zero --- and that code is \emph{attenuation-dependent}: each
\texttt{Atten} setting switches a different analog gain in front of the converter.
So the driver calibrates every unit at \emph{every} attenuation once, stores the
resulting code per \texttt{(unit, attenuation)}, and programs the matching one into
the SAR immediately before each \texttt{Read}. Calibrating at a single attenuation
and reusing that one code at the others --- the easy trap --- skews every reading
not taken at the calibrated setting; keying the code to the attenuation the caller
actually asked for keeps all four accurate.
